\chapter{Drugs and Alcohol}

\section*{Definition and Overview}
Drugs and alcohol are psychoactive substances that alter brain function, mood, perception, and behaviour. Alcohol is the most widely used recreational drug in Australia, while other drugs range from prescribed and over-the-counter medicines to illegal substances such as cannabis, methamphetamine, and heroin. This chapter provides an overview of how drugs and alcohol affect health, the risks associated with their use, and where to find help for people who want to reduce or stop their use. Understanding the effects of these substances is the first step towards making informed decisions and supporting others who may be struggling.

\section*{Epidemiology}
Alcohol and other drug use is common in Australia, with the majority of adults drinking alcohol at some level and a substantial minority using illicit drugs each year. Alcohol contributes to a significant burden of injury, violence, and chronic disease, including liver disease, cardiovascular disease, and several cancers. Tobacco use remains a leading cause of preventable death, while illicit drug use is associated with overdose, mental health problems, and blood-borne infections such as hepatitis C. Rates of use vary by age, sex, socioeconomic status, and region, and both binge drinking and daily cannabis use are particular concerns among young people.

\section*{Aetiology and Pathophysiology}
Substance use is driven by a complex mix of genetic, psychological, social, and environmental factors. The brain's reward system, which releases dopamine in response to pleasurable activities, is central to the development of dependence: drugs and alcohol hijack this system, producing intense reward and teaching the brain to seek the substance repeatedly.
\begin{itemize}
    \item \textbf{Alcohol:} a central nervous system depressant that slows brain activity, impairs judgement and coordination at low doses, and at high doses can cause blackout, respiratory depression, and coma
    \item \textbf{Stimulants:} drugs such as methamphetamine and cocaine increase alertness, energy, and euphoria but raise heart rate and blood pressure, and can cause agitation, psychosis, and cardiovascular collapse
    \item \textbf{Depressants:} benzodiazepines and opioids calm the nervous system, relieving anxiety or pain, but carry high risks of dependence, tolerance, and fatal overdose, especially when combined with alcohol
    \item \textbf{Cannabis:} produces relaxation and altered perception but can worsen anxiety and paranoia, and regular use in young people is linked to poorer educational outcomes and increased risk of psychosis
    \item \textbf{Tolerance and dependence:} repeated use leads to tolerance, so more of the drug is needed for the same effect, and withdrawal symptoms occur when use stops
    \item \textbf{Social and environmental factors:} peer influence, family history, trauma, mental illness, poverty, and availability of substances all shape patterns of use
\end{itemize}

\section*{Clinical Features}
The signs of drug and alcohol problems vary widely depending on the substance, the amount used, and the individual.
\begin{itemize}
    \item Intoxication: slurred speech, unsteadiness, drowsiness, agitation, euphoria, or aggression
    \item Withdrawal: anxiety, sweating, tremor, nausea, insomnia, and in severe cases seizures or delirium
    \item Tolerance: needing increasingly large amounts to achieve the same effect
    \item Loss of control: unsuccessful attempts to cut down, and use despite clear harm to health, work, or relationships
    \item Neglect of responsibilities and reduced performance at work, school, or home
    \item Physical health problems: liver disease, pancreatitis, heart problems, infections from injecting, and injuries sustained while intoxicated
    \item Mental health problems: depression, anxiety, paranoia, and psychosis
    \item Financial and legal problems, including drink-driving and drug offences
\end{itemize}

\section*{Differential Diagnosis}
Harmful substance use can be confused with, or masked by, several other conditions.
\begin{itemize}
    \item Mental health disorders, including depression, anxiety, bipolar disorder, and psychosis, which frequently coexist with substance use
    \item Chronic fatigue, sleep disorders, and other medical causes of low energy or poor concentration
    \item Neurological conditions that cause confusion, tremor, or unsteadiness
    \item Personality disorders and behavioural problems that may have separate causes
    \item Medication side effects that mimic the effects of intoxication or withdrawal
\end{itemize}

\section*{When to Seek Medical Care}
Anyone who is concerned about their own or another person's use of drugs or alcohol should seek advice from a GP, who can assess the situation, offer brief interventions, and arrange specialist referral. People who want to stop using opioids, benzodiazepines, or alcohol should seek medical support, because withdrawal from these substances can be dangerous.

\textbf{Call 000 (or local emergency services) for:}
\begin{itemize}
    \item Collapse, unconsciousness, or difficulty breathing after drinking or taking drugs
    \item Severe agitation, confusion, paranoia, or hallucinations
    \item Chest pain, seizures, or a suspected overdose
    \item Signs of alcohol poisoning: cold and clammy skin, slow breathing, vomiting while unconscious, or seizures
\end{itemize}

\section*{Diagnosis and Investigations}
Assessment is primarily clinical and aims to establish the pattern of use, the degree of dependence, and any associated harm.
\begin{itemize}
    \item \textbf{History:} detailed account of what substances are used, how much, how often, and the impact on health and daily life
    \item \textbf{Screening questionnaires:} tools such as the AUDIT for alcohol and the DAST for other drugs help identify risky patterns of use
    \item \textbf{Physical examination:} signs of intoxication or withdrawal, injecting sites, and evidence of liver or heart disease
    \item \textbf{Blood tests:} liver function tests, and tests for blood-borne viruses such as hepatitis B, hepatitis C, and HIV in people who inject drugs
    \item \textbf{Drug testing:} urine or breath testing may be used in treatment programmes and occupational settings
    \item \textbf{Mental health assessment:} to identify coexisting depression, anxiety, or psychosis
\end{itemize}

\section*{Management}
Management depends on the substance involved, the severity of dependence, and the person's goals, which may range from harm reduction to complete abstinence.
\begin{itemize}
    \item \textbf{Brief interventions:} short, structured conversations with a GP or counsellor that increase awareness of risks and motivate change
    \item \textbf{Counselling and psychological therapies:} cognitive behavioural therapy and motivational interviewing help people understand their use and build strategies for change
    \item \textbf{Withdrawal management (detoxification):} medically supervised withdrawal, particularly important for alcohol, benzodiazepines, and opioids, where withdrawal can be dangerous
    \item \textbf{Medication:} pharmacotherapy is available for alcohol dependence and opioid dependence, reducing cravings and preventing relapse
    \item \textbf{Residential rehabilitation:} structured programmes providing a drug-free environment, counselling, and life-skills training
    \item \textbf{Peer support:} groups such as Alcoholics Anonymous and SMART Recovery provide ongoing support and accountability
    \item \textbf{Harm reduction:} needle and syringe programmes, supervised injecting facilities, and naloxone availability reduce the harms of injecting drug use
    \item \textbf{Family support:} involving partners and families in treatment improves outcomes and helps rebuild relationships
\end{itemize}

\section*{Complications}
\begin{itemize}
    \item Dependence and addiction, with a chronic, relapsing course
    \item Overdose, which may be fatal, particularly with opioids or when drugs are combined
    \item Liver disease, including fatty liver, hepatitis, and cirrhosis from alcohol
    \item Cardiovascular disease, stroke, and hypertension
    \item Cancer, with alcohol and tobacco both being established carcinogens
    \item Mental health disorders, including depression, anxiety, and substance-induced psychosis
    \item Blood-borne infections such as hepatitis B, hepatitis C, and HIV from injecting
    \item Injuries, accidents, and violence while intoxicated
    \item Social harms, including relationship breakdown, unemployment, financial strain, and legal problems
\end{itemize}

\section*{Prognosis and Outcomes}
Substance use problems are chronic conditions, but recovery is achievable and common. Many people reduce or stop their use with brief interventions and support, while those with severe dependence often require a combination of treatment approaches over several years. Relapse is part of the recovery process rather than a failure, and each attempt at change builds skills and insight. The health harms of alcohol and other drug use are largely reversible when use stops, particularly in the early stages, making early help-seeking an important step towards better health.

\section*{Prevention and Self-Care}
\begin{itemize}
    \item Follow Australian guidelines: no more than 10 standard drinks per week and no more than 4 on any one day for healthy adults
    \item Have regular alcohol-free days and know the standard drink sizes
    \item Never drink and drive, and plan safe transport when drinking
    \item Avoid mixing drugs, including alcohol with prescription or over-the-counter medicines
    \item Model responsible drinking for children and talk openly about drugs and alcohol with teenagers
    \item Seek help early if use is affecting health, mood, work, or relationships
    \item Store medicines securely and dispose of unused medicines safely
\end{itemize}

\section*{Sources and Further Reading}
The following reputable sources provide further information for healthcare students and patients seeking reliable, current advice about drugs and alcohol.
\begin{itemize}
    \item Alcohol and Drug Foundation. \textit{Drug and alcohol information and resources.} https://adf.org.au/
    \item NHS. \textit{Alcohol misuse and drug addiction: treatment and support.} https://www.nhs.uk/conditions/alcohol-misuse/
    \item Mayo Clinic. \textit{Drug addiction: symptoms and causes.} https://www.mayoclinic.org/diseases-conditions/drug-addiction/
    \item MedlinePlus. \textit{Drugs and alcohol health information.} https://medlineplus.gov/druguseandaddiction.html
    \item World Health Organization. \textit{Alcohol and substance use.} https://www.who.int/health-topics/alcohol
    \item National Institute on Drug Abuse. \textit{Drug abuse and addiction resources.} https://nida.nih.gov/
\end{itemize}
